\documentclass[12pt,a4paper]{article}
\usepackage[margin=2.5cm]{geometry}
\usepackage{amsmath,amssymb,amsthm}
\usepackage{booktabs}
\usepackage{array}
\usepackage{multirow}
\usepackage{graphicx}
\usepackage{xcolor}
\usepackage{hyperref}
\usepackage{natbib}
\usepackage{setspace}
\usepackage{caption}
\usepackage{subcaption}

\onehalfspacing

\newtheorem{proposition}{Proposition}
\newtheorem{corollary}{Corollary}
\newtheorem{remark}{Remark}

\title{\textbf{Random Subset Methods for Forecast Combination:\\
Theory, Monte Carlo Evidence, and Empirical Applications}\\[0.5em]
\large Summary Report}
\author{Boris Kozyrev}
\date{June 2026}

\begin{document}
\maketitle
\tableofcontents
\newpage

% ===========================================================================
\section{Theoretical Framework}
% ===========================================================================

\subsection{Common-Loading Model}

The baseline model for forecast error covariance is the \emph{common-loading}
(compound-symmetry) structure:
\begin{equation}
\Sigma_e = q\,\iota_m\iota_m' + D, \qquad
D = \mathrm{diag}(\sigma_1^2,\ldots,\sigma_m^2),
\end{equation}
where $\iota_m$ is an $m$-vector of ones, $q \geq 0$ is a common variance
component shared by all $m$ forecasts, and $\sigma_i^2$ is the idiosyncratic
variance of forecaster $i$. Denote precision weights $\tau_i = \sigma_i^{-2}$,
$\bar\tau = m^{-1}\sum_i \tau_i$, and $A_S = \sum_{i\in S}\tau_i$ for a
subset $S$. Write $\tau_i = \bar\tau(1+\delta_i)$ with $\bar\delta=0$
and heterogeneity measure $V_\delta = m^{-1}\sum_i\delta_i^2$.

The infeasible minimum-variance combination weight is
$w^*_i = \bigl(\Sigma_e^{-1}\iota\bigr)_i\big/\bigl(\iota'\Sigma_e^{-1}\iota\bigr)$,
with population risk
$Q^{\mathrm{opt}} = 1/\bigl(\iota'\Sigma_e^{-1}\iota\bigr)$.
For $q=0$, $w^*_i = \tau_i/(m\bar\tau)$ and $Q^{\mathrm{opt}} = 1/(m\bar\tau)$.

\subsection{RSM Population Risk}

A Random Subset Method (RSM) draw: sample a size-$k$ subset $S$ uniformly at
random from $\{1,\ldots,m\}$, estimate sum-to-one constrained OLS weights
$w_{i,S} = \tau_i/A_S$ on $S$, and form the subset forecast. Average $G$
such draws to get $\hat f^{\mathrm{bag}}_{k,G}$.

\paragraph{Average-subset risk.}
\begin{equation}\label{eq:Qsub}
Q_k^{\mathrm{sub}} = q + \mathbb{E}_S\!\left[\frac{1}{A_S}\right].
\end{equation}

\paragraph{Infinite-bagged risk.}
Let $\bar v_{i,k} = \mathbb{E}_S[w_{i,S}\,\mathbf{1}(i\in S)]$ be the
average weight placed on forecaster $i$. Then
\begin{equation}\label{eq:Qbag}
Q_k^{\mathrm{bag},\infty} = \bar v_k'\Sigma_e\,\bar v_k.
\end{equation}

\paragraph{Weight-deviation representation.}
The excess risk of the infinite-bagged estimator over the infeasible optimum
is a quadratic form in the deviation of the average RSM weight from the
optimal weight:
\begin{equation}\label{eq:Rbag}
Q_k^{\mathrm{bag},\infty} - Q^{\mathrm{opt}}
= (\bar v_k - w^*)'\Sigma_e\,(\bar v_k - w^*) \;\geq\; 0.
\end{equation}
This representation is key: \emph{the functional form of the excess risk
is determined entirely by how the average RSM weight $\bar v_k$ deviates
from $w^*$.} The deviation vanishes as $k \to m$ (full pool), and its
dependence on $k$ drives the optimal subset size.

\paragraph{Finite-bagging bridge.}
For $G$ draws the expected loss is
\begin{equation}\label{eq:bridge}
Q_k^{\mathrm{bag},G}
= Q_k^{\mathrm{bag},\infty}
+ \frac{1}{G}\!\left(Q_k^{\mathrm{sub}} - Q_k^{\mathrm{bag},\infty}\right).
\end{equation}
This identity is exact (no approximation), and was verified in simulation to
machine precision (max error $< 9 \times 10^{-15}$). In particular,
$Q_k^{\mathrm{bag},\infty} \leq Q_k^{\mathrm{sub}}$ for all $k$: infinite
bagging is weakly better than a single subset. The gain from bagging grows
with $G$.

\subsection{Average-Subset Risk: First-Order Expansion and the Square-Root Rule}

The average-subset estimator uses a single random subset; no averaging takes
place. Its excess risk comes from $\mathbb{E}_S[1/A_S]$ in
\eqref{eq:Qsub}. Under homogeneous precision ($\tau_i = \bar\tau$), $A_S
= k\bar\tau$ exactly, so $\mathbb{E}_S[1/A_S] = 1/(k\bar\tau)$ and
\begin{equation}\label{eq:Qsub_leading}
Q_k^{\mathrm{sub}} - Q^{\mathrm{opt}}
\;\approx\; \frac{1}{k\bar\tau} - \frac{1}{m\bar\tau}
= \frac{z_k}{\bar\tau}, \qquad z_k \equiv \frac{1}{k} - \frac{1}{m}.
\end{equation}
The excess risk is \emph{linear} in $z_k$ (effectively $\alpha=1$).

Including finite-sample weight-estimation inflation $(T-1)/(T-k)$, the
expected in-sample loss is
\begin{equation}\label{eq:Lsub}
L_k^{\mathrm{sub}}
\;\approx\; \left(Q^{\mathrm{opt}} + \frac{z_k}{\bar\tau}\right)
            \frac{T-1}{T-k}.
\end{equation}
Setting $\partial L_k^{\mathrm{sub}}/\partial k = 0$ and approximating
$z_k \approx 1/k$ for $k \ll m$ yields
\begin{equation}\label{eq:kstar_sub}
k^*_{\mathrm{sub}} \;\approx\;
\left(\frac{(T-1)}{\bar\tau\,Q^{\mathrm{opt}}}\right)^{1/2}
\;\propto\; T^{1/2}.
\end{equation}
The average-subset estimator has a \emph{square-root} optimal subset size.

\subsection{Infinite-Bagged Risk under Small Heterogeneity: the Cube-Root Rule}

Bagging replaces the random single-subset weight $w_{i,S}$ with its
expectation $\bar v_{i,k}$, concentrating the weight and reducing variance.
To find the functional form of the excess risk \eqref{eq:Rbag}, we compute
the weight deviation $\Delta_{i,k} = \bar v_{i,k} - w^*_i$.

\paragraph{Weight deviation under small heterogeneity.}
With $\tau_i = \bar\tau(1+\delta_i)$ and $V_\delta \ll 1$, a first-order
Taylor expansion of $\tau_i/A_S$ around the homogeneous case gives
\begin{equation}\label{eq:delta}
\Delta_{i,k} = \bar v_{i,k} - w^*_i
\;\approx\; -\frac{\delta_i\,z_k}{m\bar\tau}.
\end{equation}
The deviation is proportional to $z_k$: it vanishes as $k \to m$ (the
bagged estimator recovers the optimum when all forecasters are included) and
is scaled by the idiosyncratic heterogeneity $\delta_i$.

\paragraph{Excess risk as a quadratic in $z_k$.}
Substituting \eqref{eq:delta} into \eqref{eq:Rbag} and using
$\Sigma_e = q\iota\iota' + D$:
\begin{equation}\label{eq:Rbag_smallhet}
Q_k^{\mathrm{bag},\infty} - Q^{\mathrm{opt}}
\;\approx\; \Delta'D\,\Delta
= \frac{z_k^2}{m^2\bar\tau^2}\sum_i \sigma_i^2\delta_i^2
\;=\; C_\delta\!\left(\frac{1}{k}-\frac{1}{m}\right)^{\!2}.
\end{equation}
The common-loading term $q(\iota'\Delta)^2$ vanishes exactly because
$\sum_i\delta_i = 0$. The coefficient $C_\delta = V_\delta/(m\bar\tau^3)$
captures the dispersion of precision weights.

The excess risk is \emph{quadratic} in $z_k$ (i.e.\ $\alpha = 2$ in a
log-log sense), in contrast to the linear $\alpha = 1$ decay for $Q^{\mathrm{sub}}$.
Bagging effectively \emph{squares} the decay rate of excess risk with $k$.

\paragraph{Optimal $k$ for the bagged estimator.}
The total loss of the infinite-bagged estimator, including finite-sample
inflation, is
\begin{equation}\label{eq:Lbag}
L_k^{\mathrm{bag}} \;\approx\;
\left[Q^{\mathrm{opt}} + C_\delta\!\left(\frac{1}{k}-\frac{1}{m}\right)^{\!2}\right]
\frac{T-1}{T-k}.
\end{equation}
Setting $\partial L_k^{\mathrm{bag}}/\partial k = 0$ and approximating
$z_k \approx 1/k$ for $k \ll m$ yields the \emph{cube-root optimal subset size}:
\begin{equation}\label{eq:kstar_bag}
k^*_{\mathrm{bag}} \;\approx\; \left(\frac{2C_\delta\,T}{Q^{\mathrm{opt}}}\right)^{1/3}
\;\propto\; T^{1/3}.
\end{equation}
The $T$-scaling exponent drops from $\tfrac{1}{2}$ (single subset) to
$\tfrac{1}{3}$ (infinite bag): bagging makes larger subsets optimal relative
to the average-subset rule.

\subsection{General Heterogeneity and the Power-Law Approximation}

When heterogeneity is large, equation \eqref{eq:Rbag_smallhet} no longer
holds and the weight deviations $\Delta_{i,k}$ do not scale as $\delta_i z_k$.
Empirically, the log-log regression
\begin{equation}\label{eq:loglog}
\log\!\left(Q_k^{\mathrm{bag},\infty} - Q^{\mathrm{opt}}\right)
\;\approx\; \hat\alpha \log z_k + \mathrm{const}
\end{equation}
yields an \emph{effective log-log slope} $\hat\alpha$ that lies in $[1,2]$
depending on the degree of heterogeneity.
Small heterogeneity gives $\hat\alpha \approx 2$ (the small-het formula,
eq.~\eqref{eq:Rbag_smallhet}); high heterogeneity produces $\hat\alpha$
closer to 1.

\begin{remark}[Effective vs.\ exact exponent]
\label{rem:alpha}
$\hat\alpha$ is an OLS slope of \eqref{eq:loglog} over the chosen $k$-grid,
not an exact mathematical exponent. In the dominant-forecaster limit
($\tau_1 \gg \tau_j$ for $j\neq 1$), the exact curve is
$R_k = C'\bigl(1-k/m\bigr)^2$,
which is quadratic in $(1-k/m)$ but not a pure power of $z_k = (m-k)/(km)$.
Near $k=m$, it behaves locally as $C'm^2z_k^2$, but an OLS regression over
the full grid gives a slope near 1 because $(1-k/m)^2$ and $z_k$ have
different curvatures over the empirically relevant range of $k$.
We therefore call $\hat\alpha$ the \emph{effective} slope; the bounds
$\hat\alpha \in [1,2]$ describe the OLS projection, not pointwise power-law
identities. Simulations consistently recover $\hat\alpha \in [1,2]$ across
all heterogeneity levels studied, with $\hat\alpha$ increasing from $\approx1$
at low $H$ to $\approx1.7$ at high $H$ (Table~\ref{tab:exo_main}).%
\footnote{The dominant-forecaster weight deviation for $j\neq 1$ is
$\bar{v}_{j,k} = P(j\in S)\cdot P(1\notin S\mid j\in S)\cdot(1/k)
= (k/m)\cdot(m-k)/(m-1)\cdot(1/k) = (m-k)/(m(m-1))$,
so $\Delta_{j,k} = \bar{v}_{j,k} - w_j^* \approx (m-k)/(m(m-1))$ since
$w_j^*\to 0$ as $\tau_1\to\infty$.
The resulting $R_k \approx (m-1)\Delta_{j,k}^2/\tau_j = (m-k)^2/(m^2(m-1)\tau_j)$
is \emph{not} a power of $z_k$ globally; locally near $k=m$ it is $\propto z_k^2$.}
\end{remark}

Minimising the generalised loss $L_k \approx (A + Cz_k^{\hat\alpha})(T-1)/(T-k)$
using the effective slope $\hat\alpha$ gives
\begin{equation}\label{eq:kstar_gen}
k^* \;\approx\; \left(\frac{C\hat\alpha\,T}{A}\right)^{1/(\hat\alpha+1)},
\end{equation}
so the effective $T$-scaling exponent $1/(\hat\alpha+1)$ ranges from
$\tfrac{1}{3}$ (small-het, $\hat\alpha=2$) to $\tfrac{1}{2}$ (strongly
heterogeneous, $\hat\alpha=1$). The cube-root rule \eqref{eq:kstar_bag}
is the \emph{lower bound} on the effective scaling exponent.

\begin{remark}[Effective scaling caveat]
Formula~\eqref{eq:kstar_gen} uses the fitted power-law approximation and
should be read as an effective scaling rule. It accurately describes the
optimal subset size when the log-log fit is tight (as confirmed empirically
for all four DGPs in Table~\ref{tab:exo_main} with $R^2>0.97$). Under the
exact dominant-forecaster curve the scaling exponent is approximately
$1/(\hat\alpha+1)$ with the effective $\hat\alpha$, giving the same
practical guidance.
\end{remark}

In both cases $k^*$ is increasing in $T$ and in $\hat\alpha$ (higher
heterogeneity requires a larger subset to cover the precision-weight
distribution).

% ===========================================================================
\section{Exogenous Monte Carlo (Known $\Sigma_e$)}
% ===========================================================================

\subsection{Design}

The exogenous experiment takes $\Sigma_e$ as given (no estimation). The DGP
follows the common-loading structure with $m = 50$, $T = 240$ (baseline).
Forecast errors are generated as
\begin{equation}
x_{it} = B_{\mathrm{scale}}\cdot F_t + H_{\mathrm{scale}} \cdot \varepsilon_{it},
\end{equation}
where $F_t \sim N(0,1)$ is a common factor and $\varepsilon_{it}$ are
idiosyncratic errors. The parameter $B \in \{3,8\}$ controls the common
variance $q$ and $H \in \{1.05, 6\}$ controls the spread of individual
precisions $\tau_i$ (heterogeneity). Simulations use $N_{\mathrm{rep}} = 500$
replications and $G = 800$ bags. Key theory quantities ($Q^{\mathrm{opt}}$,
$C_\delta$, $k^*$ formulae) are computed analytically from the true $\Sigma_e$.

\subsection{Main Results}

Table~\ref{tab:exo_main} summarises the four DGP combinations, using all
available population-level and finite-sample statistics. The column
``$k^*$(grid)'' is the minimiser of $Q_k^{\mathrm{bag},\infty}$ over
$k \in \{1,\ldots,m\}$; ``$k^*$(emp)'' is the minimiser of the finite-sample
MSPE over 500 replications; ``Gain'' is the percentage improvement of the
optimal RSM over the equal-weighted mean.

\begin{table}[h]
\centering
\caption{Exogenous MC: $m=50$, $T=240$, 500 replications, 800 bags.}
\label{tab:exo_main}
\begin{tabular}{llcccccc}
\toprule
$B$ & $H$ & $k^*$(grid) & $k^*$(emp) & Gain vs EW
    & $\hat\alpha$ & $R^2$ & Approx.\ err \\
\midrule
3 & 1.05 & 1 & 3 & $0.1\%$ & 1.12 & 0.98 & 90\% \\
8 & 1.05 & 1 & 3 & $0.1\%$ & 1.10 & 0.98 & 90\% \\
3 & 6    & 3 & 6 & $12.7\%$ & 1.72 & 0.98 & 75\% \\
8 & $3.5^*$ & 4 & 8 & $14.9\%$ & 1.70 & 0.99 & 77\% \\
\bottomrule
\multicolumn{8}{l}{\footnotesize $^*$ Target $H=6$; parameterisation achieved $H_{\mathrm{actual}}=3.51$.}\\
\multicolumn{8}{l}{\footnotesize $\hat\alpha$: OLS slope in $\log(Q_k^{\mathrm{bag},\infty}-Q^{\mathrm{opt}}) \sim \hat\alpha\log z_k$;
  Approx.\ err: mean $|R_k - C_\delta z_k^2|/R_k$ over $k$.}
\end{tabular}
\end{table}

\paragraph{Theory alignment.}
Several predictions are confirmed:
\begin{itemize}
\item \textbf{Finite-bagging bridge}: equation \eqref{eq:bridge} holds to
  machine precision ($<9\times10^{-15}$) for all $k$ and all DGPs.
\item \textbf{Cube-root formula}: $k^*(\text{formula}) = k^*(\text{grid})$
  in all four DGPs, confirming that \eqref{eq:kstar_bag} locates the correct
  minimiser even when the small-het approximation fails by a factor of
  5--100$\times$.
\item \textbf{Power-law}: $\hat\alpha \in [1,2]$ with $R^2 > 0.97$ in all
  cases. $\hat\alpha$ increases with $H$ (from $\approx 1.1$ at $H\approx1$
  to $\approx 1.7$ at $H\approx6$), consistent with the theoretical bounds.
  It depends on $H$ (heterogeneity) but \emph{not} on $B$ (common factor).
\item \textbf{Gains grow with heterogeneity}: 0.1\% at $H=1.05$ vs
  12--15\% at $H=3.5$--6. RSM is most valuable when forecasters differ
  substantially in quality.
\end{itemize}

\paragraph{Small-het approximation failure.}
The approximation error (90\% at $H=1.05$, 75--77\% at $H=3.5$--6) is large.
The true excess risk exceeds $C_\delta z_k^2$ substantially at moderate $k$,
as the Taylor expansion of $\mathbb{E}_S[1/A_S^2]$ involves higher-order
heterogeneity moments. Despite this, the cube-root formula gives the exact
minimiser because the leading bias term in the approximation is constant in
$k$ and does not affect the location of the minimum.

\subsection{$T$-Scaling and the Cube-Root Law}

The $T$-scaling experiment varies $T \in \{80, 240, 640, 1000\}$ at
$B=3$, $H=1.05$. Table~\ref{tab:exo_scaling} shows the continuous cube-root
formula value, the grid-search optimal, and the empirical optimum.

\begin{table}[h]
\centering
\caption{Exogenous $T$-scaling: $B=3$, $H\approx1.05$, $m=50$.}
\label{tab:exo_scaling}
\begin{tabular}{lcccc}
\toprule
$T$ & $k^*$(formula, cont.) & $k^*$(grid, bag$^\infty$) & $k^*$(emp) & Gain \\
\midrule
80   & 0.57 & 1 & 2 & $0.02\%$ \\
240  & 0.82 & 1 & 2 & $0.02\%$ \\
640  & 1.13 & 1 & 4 & $0.06\%$ \\
1000 & 1.31 & 1 & 3 & $0.03\%$ \\
\bottomrule
\end{tabular}
\end{table}

At $H=1.05$, the population excess risk is virtually zero for any $k$
(gains $< 0.1\%$), so the loss function is essentially flat and the
population optimum is $k^*=1$ throughout. The cube-root formula
correctly predicts this: the continuous value ranges from 0.57 to 1.31,
confirming that $k^*=1$ is optimal at every $T$ studied. The cube-root
growth rate is present but invisible because the signal (heterogeneity)
is negligible.

The empirical $k^* = 2$--4 reflects simulation noise on a near-flat loss
surface, not a structural deviation from the theory. Detecting the cube-root
law requires sufficient heterogeneity ($H \gg 1$); the next subsection
confirms this at $H=6$.

\subsection{$T$-Scaling at $H=6$: The Cube-Root Law Detected}
\label{sec:h6scaling}

To confirm that the cube-root scaling law is present but masked at low
heterogeneity, we repeat the $T$-scaling experiment at $H=6$ ($B=3$,
$m=50$, 500 replications). Table~\ref{tab:exo_scaling_h6} reports
the results.

\begin{table}[h]
\centering
\caption{Exogenous $T$-scaling: $B=3$, $H\approx6.08$, $m=50$, 500 reps.}
\label{tab:exo_scaling_h6}
\begin{tabular}{lcccc}
\toprule
$T$ & $k^*$(formula, cont.) & $k^*$(grid, bag$^\infty$) & $k^*$(emp) & Gain \\
\midrule
80   & 2.0 & 2 & 4  & 12.4\% \\
240  & 2.9 & 3 & 6  & 12.1\% \\
640  & 4.0 & 4 & 8  & 12.8\% \\
1000 & 4.6 & 5 & 10 & 12.8\% \\
\bottomrule
\end{tabular}
\end{table}

\noindent
The log-log regression of $k^*(\text{grid})$ on $T$ gives slope
$\hat\beta = 0.352$ ($R^2 = 0.994$) --- essentially exactly $\tfrac{1}{3}$.
The result for $k^*(\text{emp})$ is identical ($\hat\beta = 0.352$,
$R^2 = 0.994$). In contrast to the $H=1.05$ case (where gains are
$<0.1\%$ and $k^*=1$ throughout), at $H=6$ gains are approximately
constant at 12--13\% across all $T$, and the cube-root scaling is clearly
detectable. This confirms that the law is structural, not an artefact of
heterogeneity level: it holds whenever heterogeneity is non-trivial.

\paragraph{Connection to the formula.}
The continuous cube-root formula \eqref{eq:kstar_bag} predicts
$k^* \in \{2.0, 2.9, 4.0, 4.6\}$ for $T \in \{80, 240, 640, 1000\}$,
matching the grid search exactly after rounding. The effective $\hat\alpha$
at $H=6$ is 1.72 (Table~\ref{tab:exo_main}), giving a predicted scaling
exponent of $1/(1.72+1) = 0.368$, consistent with the empirical 0.352.

% ===========================================================================
\section{Empirical Application: FRED-MD GDP Nowcasting}
% ===========================================================================

\subsection{Data and Setup}

We use 102 real-time quarterly vintages of FRED-MD (September 1999 --
March 2025), excluding Q4 2003 (missing data). At each vintage, the
estimation sample covers approximately 1970Q1 to the forecast origin,
giving $T_{\mathrm{eff}} \approx 79$--183 quarters (median $\approx 100$).
The predictor pool contains $m \approx 126$ monthly FRED-MD series,
aggregated to quarterly frequency.

\paragraph{Individual forecasts (bivariate).}
For each predictor series $j$ and time $t$, we estimate
\begin{equation}
y_t = \alpha_j + \beta_j x_{j,t} + \varepsilon_t,
\end{equation}
where $y_t$ is log-differenced real GDP growth. This yields $m$ individual
forecasts per period.

\paragraph{Individual forecasts (AR+X).}
Following \citet{ElliottEtAl2015}, we also produce forecasts from
\begin{equation}
y_t = \alpha_j + \varphi_1 y_{t-1} + \cdots + \varphi_p y_{t-p}
      + \beta_j x_{j,t} + \varepsilon_t,
\end{equation}
where $p$ is chosen by AIC at each $t$.

\paragraph{RSM combination.}
At each vintage and period $t$, RSM samples $k$ individual forecasts and
runs sum-to-one constrained OLS. We average $B=1000$ draws. For robustness
we sweep $k \in \{2,\ldots,55\}$ with $B=200$ draws.

\paragraph{Benchmarks.} Equal-weighted mean, Lasso, Ridge, Random Forest.

\subsection{Results: Bivariate RSM}

\begin{table}[h]
\centering
\caption{FRED-MD GDP: bivariate RSM, full sample (102 quarters).}
\label{tab:gdp_biv_full}
\begin{tabular}{lcccc}
\toprule
Model & MAFE & Sig. & RMSFE & Sig. \\
\midrule
Mean           & 0.54 &     & 1.10 &    \\
RSM            & 0.45 & **  & 0.76 & *  \\
Lasso          & 0.51 &     & 0.73 & *  \\
Ridge          & 0.49 & *   & 0.98 &    \\
Random Forest  & 0.55 &     & 1.24 &    \\
\bottomrule
\multicolumn{5}{l}{\footnotesize $k = \lfloor\sqrt{T}\rfloor \approx 9$.
  *$p<0.10$, **$p<0.05$ (DM test vs.\ Mean).}
\end{tabular}
\end{table}

\begin{table}[h]
\centering
\caption{FRED-MD GDP: bivariate RSM, excluding COVID-19 (Q2--Q3 2020).}
\label{tab:gdp_biv_nocov}
\begin{tabular}{lcccc}
\toprule
Model & MAFE & Sig. & RMSFE & Sig. \\
\midrule
Mean           & 0.42 &     & 0.60 &     \\
RSM            & 0.38 & **  & 0.53 & **  \\
Lasso          & 0.45 &     & 0.58 &     \\
Ridge          & 0.39 & **  & 0.54 & **  \\
Random Forest  & 0.41 &     & 0.57 & *   \\
\bottomrule
\end{tabular}
\end{table}

\paragraph{RMSFE by $k$.}
The RMSFE$(k)$ curve is U-shaped. Key values:
\begin{itemize}
\item Full sample: $k^*(\text{RMSFE}) = 19$, min RMSFE $= 0.662$.
      At $k=\sqrt{T} \approx 9$: RMSFE $= 0.762$.
\item No-COVID: $k^*(\text{RMSFE}) = 21$, min RMSFE $= 0.499$.
      At $k = 9$: RMSFE $= 0.529$.
\end{itemize}

\subsection{Results: AR+X RSM}

\begin{table}[h]
\centering
\caption{FRED-MD GDP: AR+X RSM, full sample.}
\label{tab:gdp_arx_full}
\begin{tabular}{lcccc}
\toprule
Model & MAFE & Sig. & RMSFE & Sig. \\
\midrule
Mean           & 0.56 &     & 1.27 &     \\
RSM (AR+X)     & 0.45 & **  & 0.76 & *   \\
Lasso          & 0.52 &     & 0.77 &     \\
Ridge          & 0.50 & **  & 1.02 & *   \\
Random Forest  & 0.57 &     & 1.40 &     \\
\bottomrule
\end{tabular}
\end{table}

\begin{table}[h]
\centering
\caption{FRED-MD GDP: AR+X RSM, excluding COVID-19.}
\label{tab:gdp_arx_nocov}
\begin{tabular}{lcccc}
\toprule
Model & MAFE & Sig. & RMSFE & Sig. \\
\midrule
Mean           & 0.41 &     & 0.57 &     \\
RSM (AR+X)     & 0.39 & *   & 0.52 & *   \\
Lasso          & 0.47 &     & 0.63 &     \\
Ridge          & 0.39 & **  & 0.54 & **  \\
Random Forest  & 0.40 &     & 0.57 &     \\
\bottomrule
\end{tabular}
\end{table}

\paragraph{Bivariate vs.\ AR+X comparison.}
\begin{center}
\begin{tabular}{lcccc}
\toprule
 & \multicolumn{2}{c}{Bivariate} & \multicolumn{2}{c}{AR+X} \\
\cmidrule(lr){2-3}\cmidrule(lr){4-5}
Sample & $k^*$ & min RMSFE & $k^*$ & min RMSFE \\
\midrule
Full    & 19 & 0.662 & 19 & \textbf{0.617} \\
No-COVID & 21 & \textbf{0.499} & 16 & 0.506 \\
\bottomrule
\end{tabular}
\end{center}

Adding AR lags lowers the full-sample minimum RMSFE by 6.8\%. In the
no-COVID sample, the benefit disappears (AR lags add estimation noise
without large cyclical episodes to track). The AR+X $k^*=16$ (no-COVID)
is closer to the MC range of 11--14 than the bivariate $k^*=21$.

\subsection{Connection to the Monte Carlo Evidence}

\begin{table}[h]
\centering
\caption{Comparison of $k^*(\text{RMSFE})$ across MC and empirical settings.}
\label{tab:kstar_compare}
\begin{tabular}{lcc}
\toprule
Setting & $T$ & $k^*(\text{RMSFE})$ \\
\midrule
GDP bivariate, full sample  & $\approx 100$ & 19 \\
GDP bivariate, no-COVID     & $\approx 100$ & 21 \\
GDP AR+X, full sample       & $\approx 100$ & 19 \\
GDP AR+X, no-COVID          & $\approx 100$ & 16 \\
\bottomrule
\end{tabular}
\end{table}

The empirical $k^*$ (16--21 at $T \approx 100$) is larger than what the
exogenous population theory predicts ($k^* \approx 1$--5 for $T=240$),
for two structural reasons:

\begin{enumerate}
\item \textbf{Underpowered covariance estimation}: GDP has
  $m \approx 126$ predictors with $T \approx 100$ observations
  ($m/T \approx 1.3$). Estimating combination weights is much harder
  than in the exogenous MC (where $\Sigma_e$ is known). The additional
  estimation noise pushes $k^*$ substantially rightward.
\item \textbf{Individual forecast noise}: each predictor's individual
  forecast carries its own OLS estimation error (from the
  bivariate/AR+X regression). Lower signal-to-noise per individual
  forecast makes larger subsets more attractive.
\end{enumerate}

The AR+X specification reduces but does not eliminate this gap: it
improves individual forecast quality, which brings the no-COVID $k^*$
from 21 to 16. Both remain well above the exogenous population optimum,
which is characteristic of high-dimensional, low-$T$ empirical settings.

% ===========================================================================
\section{SPF: The Near-Optimal Case}
% ===========================================================================

\subsection{Setup and Motivation}

The theory predicts that RSM gains are negligible when forecasters are
near-homogeneous: at $H \approx 1$, the cube-root formula gives $k^* \approx 1$
and gains are $<0.1\%$ (Table~\ref{tab:exo_main}). The Survey of Professional
Forecasters (SPF) provides a direct test of this prediction: SPF respondents
are trained professional economists using similar information sets, making
near-homogeneity ($H \approx 1$) the natural prior.

We apply RSM directly to SPF responses, replicating the setup of
\citet{ConfittiEtAl2015}. The sample covers Q1 1985 -- Q4 2024, with
$N_t \approx 30$--53 respondents per round (median 36). Each forecaster's
own submission is the individual forecast (no regression estimation).
$B=1000$ draws per round. Coverage constraint: $k \leq k_{\max}$ where
$k_{\max}$ is the largest $k$ at which $\geq 80\%$ of rounds have $N_t \geq k$.

Four target variables: real GDP growth ($h=1$ and $h=4$ quarters ahead)
and CPI inflation ($h=1$ and $h=4$ quarters ahead).

\subsection{Results: Hard to Beat the Mean}

\begin{table}[h]
\centering
\caption{SPF RSM vs.\ equal-weighted mean; OOS 1985Q1--2024Q4.}
\label{tab:spf}
\begin{tabular}{llccccc}
\toprule
Series & $h$ & Median $N$ & $k^*$ & $k=\!\sqrt{N}$ & RMSFE (Mean) & RMSFE (RSM) \\
\midrule
RGDP & 1 & 36 & 3  & 6 & 6.405 & 6.392 ($-0.2\%$) \\
RGDP & 4 & 36 & 13 & 6 & 2.208 & 2.205 ($-0.1\%$) \\
CPI  & 1 & 35 & 10 & 6 & 0.974 & 0.969 ($-0.5\%$) \\
CPI  & 4 & 35 &  3 & 6 & 1.471 & 1.470 ($-0.1\%$) \\
\bottomrule
\multicolumn{7}{l}{\footnotesize No DM tests significant at 10\% level.}
\end{tabular}
\end{table}

Gains are uniformly below 0.5\% and insignificant. This confirms the
main theoretical prediction: \emph{it is hard to beat the equal-weighted
mean when forecasters are near-homogeneous.} Several features of the SPF
environment explain the small gains:

\begin{enumerate}
\item \textbf{Near-homogeneous precision.} Professional forecasters use
  similar models and information sets; heterogeneity $H \approx 1$. The
  MC shows that at $H=1.05$ gains are $<0.1\%$, matching the SPF results.

\item \textbf{No regression estimation noise.} SPF individual forecasts
  are direct submissions, without additional estimation error. In the
  FRED-MD setting, each individual forecast carries OLS estimation
  noise that RSM partially averages out. Removing that noise removes a
  channel through which RSM gains arise.

\item \textbf{Small pool.} At $N \approx 36$ respondents, distinct-subset
  diversity is limited. The combinatorial benefit of RSM grows with $m$;
  at $m=36$ it is much smaller than at $m=126$ (FRED-MD) or $m=50$ (MC).
\end{enumerate}

Despite the small magnitude, the SPF results confirm the \emph{direction}
of all MC predictions: RSM never underperforms the mean; $k^*$ is interior
(not at the boundary); and the $\sqrt{N}$ rule of thumb overestimates $k^*$
in the low-heterogeneity cases (RGDP $h=1$, CPI $h=4$) and underestimates it
in the moderate-heterogeneity case (CPI $h=1$).

% ===========================================================================
\section{Overall Alignment: Theory, MC, and Empirics}
% ===========================================================================

Table~\ref{tab:alignment} provides a unified scorecard of the theoretical
predictions across the three empirical settings.

\begin{table}[h]
\centering
\caption{Theory--MC--Empirics alignment.}
\label{tab:alignment}
\small
\begin{tabular}{p{5.8cm}ccc}
\toprule
Prediction & Exog.\ MC & Empirical (GDP) & SPF \\
\midrule
Finite-bagging bridge (exact)              & \checkmark & n/a  & n/a  \\
Cube-root formula $= k^*$(grid)            & \checkmark & $\times$ & n/a \\
Power-law $\hat\alpha \in [1,2]$           & \checkmark & n/a  & n/a  \\
$\hat\alpha$ depends on $H$, not $B$       & \checkmark & n/a  & n/a  \\
$k^*$ scales as $T^{1/3}$ at $H=6$         & \checkmark & n/a  & n/a \\
$k^*$ scales as $T^{1/(\hat\alpha+1)}$     & \checkmark & $\sim$   & n/a \\
RSM beats equal weights                    & \checkmark & \checkmark & \checkmark \\
Near-homogeneous: gains $<1\%$             & \checkmark & n/a  & \checkmark \\
$\sqrt{T}$ rule of thumb for $k^*$         & $\sim$     & $\sim$   & $\sim$ \\
Gains grow with heterogeneity              & \checkmark & \checkmark & \checkmark \\
U-shaped RMSFE$(k)$ in practice            & n/a  & \checkmark & n/a \\
\bottomrule
\multicolumn{4}{l}{\footnotesize \checkmark: confirmed; $\sim$: partial;
  $\times$: not confirmed; n/a: not applicable.}
\end{tabular}
\end{table}

\paragraph{Summary.}
The exogenous MC validates the core theoretical machinery: the exact
finite-bagging bridge, the cube-root formula, and the power-law decay
of excess risk all hold precisely. The key empirical finding is
that heterogeneity ($H$) is the decisive parameter: at $H \approx 1$
(SPF, near-homogeneous forecasters) RSM gains are negligible and it is
hard to beat the simple mean. At $H \gg 1$ (FRED-MD GDP, heterogeneous
predictors) RSM produces statistically significant gains of 10--30\%.

The FRED-MD $k^*$ (16--21) exceeds the exogenous population prediction
because estimation of combination weights is expensive when $m/T > 1$:
the optimal $k$ must balance approximation bias against the additional
noise from fitting constrained OLS on each subset. The AR+X specification
brings $k^*$ closer to the MC range by improving individual forecast quality,
confirming the structural interpretation.

\appendix

\section{Data Sources}
\begin{itemize}
\item \textbf{FRED-MD}: Real-time monthly vintages, September 1999 -- March
  2025. 103 vintages; Q4 2003 excluded (missing data). McCracken and Ng (2016).
\item \textbf{FRED-QD}: Quarterly vintages matched to monthly by vintage date.
\item \textbf{SPF}: Philadelphia Fed Survey of Professional Forecasters.
  RGDP file (columns RGDP1--RGDP5) and CPI file (columns CPI1--CPI5).
  Actuals: GDPC1 and CPIAUCSL from FRED (downloaded June 2026).
\end{itemize}

\section{$m$-Scaling: Pool Size and Optimal $k^*$}
\label{app:mscaling}

The cube-root formula \eqref{eq:kstar_bag} scales with $T$ but does not
explicitly depend on $m$. In practice, larger pools require larger
subsets because (i) the precision weight distribution is more spread out
and (ii) estimating combination weights is noisier when $m$ is large
relative to $T$. To quantify the $m$-dependence, we run the endogenous MC
at $T=200$ fixed and $m \in \{20, 30, 50, 80, 120\}$ with paper-level
settings ($N_{\mathrm{rep}}=500$, $K_{\mathrm{RS}}=500$, OOS length $=50$,
B4 coefficients).

\begin{table}[h]
\centering
\caption{$m$-scaling: $k^*(\mathrm{RMSFE})$ by pool size, $T=200$ fixed.}
\label{tab:m_scaling}
\begin{tabular}{lcccccc}
\toprule
$\Sigma_e$ & $m=20$ & $m=30$ & $m=50$ & $m=80$ & $m=120$ & log-log slope \\
\midrule
Toeplitz  & 8 & 8 & 12 & 17 & 22 & $0.611$ ($R^2=0.96$, $p=0.004$) \\
Factor    & 8 & 10 & 14 & 20 & 18$^*$ & $0.512$ ($R^2=0.91$, $p=0.012$) \\
\bottomrule
\multicolumn{7}{l}{\footnotesize $^*$Factor $m=120$: $k^*$ dips to 18 from 20 at $m=80$,
  likely grid-step noise (step size 4).}\\
\multicolumn{7}{l}{\footnotesize Log-log slope: OLS of $\log k^*$ on $\log m$.
  Gains stabilise at 10--11\% (Toeplitz) and 4--7\% (Factor) across $m$.}
\end{tabular}
\end{table}

\noindent
$k^*$ rises sub-linearly with $m$ for both $\Sigma_e$ types, with
log-log slopes of 0.61 (Toeplitz) and 0.51 (Factor) --- well below 1
(linear) and above 0.33 (cube-root). This confirms that larger pools
require larger subsets, with the rate between the $T^{1/3}$ and
$T^{1/2}$ predictions. The Kan-Zhou correction factor
$h(m,T)=(T-2)/(T-m-2)$ (which captures estimation noise from estimating
$\Sigma_e$ from $m$ series) rises from 1.11 at $m=20$ to 2.54 at $m=120$,
quantitatively consistent with the observed $k^*$ trajectory for Toeplitz.

\section{Software and Replication}
All code is in R ($\geq 4.2$). Key packages: \texttt{dplyr}, \texttt{glmnet},
\texttt{ranger}, \texttt{forecast}, \texttt{readxl}. All random number streams
are seeded. Scripts:
\begin{itemize}
\item Exogenous MC: \texttt{rsm\_common\_loading\_grid\_weights\ldots\_final.R}
\item GDP bivariate: \texttt{GDP\_constrained/GDP\_all\_k\_results.R}
\item GDP AR+X: \texttt{GDP\_constrained/GDP\_ar\_bivariate\_allk.R}
\item SPF: \texttt{GDP\_constrained/SPF\_RSM\_analysis.R}
\item $m$-scaling MC: \texttt{rsm\_m\_scaling.R}
\end{itemize}

\begin{thebibliography}{99}

\bibitem[Conflitti et al.(2015)]{ConfittiEtAl2015}
Conflitti, C., De Mol, C., \& Giannone, D. (2015).
Optimal combination of survey forecasts.
\textit{International Journal of Forecasting}, 31(4), 1096--1103.

\bibitem[Elliott et al.(2013)]{ElliottEtAl2015}
Elliott, G., Gargano, A., \& Timmermann, A. (2013).
Complete subset regressions.
\textit{Journal of Econometrics}, 177(2), 357--373.

\bibitem[Kan and Zhou(2007)]{KanZhou2007}
Kan, R., \& Zhou, G. (2007).
Optimal portfolio choice with parameter uncertainty.
\textit{Journal of Financial and Quantitative Analysis}, 42(3), 621--656.

\end{thebibliography}

\end{document}
